% Paste this fragment after \appendix in the Prism manuscript source.
% It uses only standard LaTeX and amsmath commands.

\section{Generation of Target-Suppressed Images $I^{-k}$}

Let $I$ denote an original image and $I^{-k}$ its counterfactual counterpart in which the visual evidence associated with target $k$ is selectively suppressed. Because the appearance, acquisition process, and spatial characteristics differ substantially across breast ultrasonography, brain MRI, chest CT, and natural-image datasets, $I^{-k}$ was generated using a domain-specific procedure rather than a single universal transformation.

\subsection{Breast ultrasonography}

For the BUS-CoT and UDIAT breast ultrasound datasets, lesion-suppressed images were generated using the \texttt{stable-diffusion-v1-5/stable-diffusion-inpainting} model in half precision. A \texttt{DDIMScheduler} was instantiated while retaining the scheduler configuration of the pretrained model. Inference was performed for 50 denoising steps with a classifier-free guidance scale of 5.0, and the safety checker was disabled. The documented test-set generation run used a batch size of two images per GPU and was divided into three independent shards. A total of 906 images were generated, consisting of 519 images from the initial run and 387 images from the resumed run.

Instead of conditioning the diffusion model on a manually written natural-language prompt, we used a fixed normal breast ultrasound image, \texttt{assets/busi\_normal\_1.png}, as the normal-reference condition. The same reference image was used for all input images. Following the image-to-prompt conversion procedure introduced in \emph{The CLIP Model Is Secretly an Image-to-Prompt Converter}, the 1,024-dimensional vision CLS embedding extracted by \texttt{openai/clip-vit-large-patch14} was mapped to the 768-dimensional Stable Diffusion text hidden space using the closed-form transformation

\[
W_{\mathrm{map}}
=
\operatorname{pinv}\!\left(W_{\mathrm{text}}\right)
W_{\mathrm{visual}}.
\]

The singular-value decomposition cutoff was set to 0.3, retaining 607 singular values. The resulting 768-dimensional image-derived token was inserted into the Stable Diffusion text-conditioning sequence. Specifically, token position 0 of the native empty prompt was retained, whereas positions 1 through 76 were filled by repeating the same normal-image-derived token. Thus, no natural-language positive prompt was used. The negative prompt was represented by the native Stable Diffusion embedding of an empty string, and the architecture and parameters of the Stable Diffusion model itself were not modified.

The original lesion mask was loaded in grayscale and dilated using \texttt{dilation\_px=9}. In this implementation, the operation corresponds to a $19 \times 19$ square maximum filter and therefore expands the mask by up to nine pixels in each direction. The source image was resized to $512 \times 512$ using Lanczos interpolation, whereas the binary mask was resized using nearest-neighbor interpolation. After inpainting, the generated image was resized back to the original spatial resolution using Lanczos interpolation. Output filenames were preserved. Because JPEG quality and chroma-subsampling parameters were not explicitly specified in the generation script, the corresponding Pillow defaults were used.

Generation was made deterministic at the image level. For an image at global position $i$ in the sorted input list, the random seed was set to

\[
s_i = 20260728 + i.
\]

Because the seed depended on the global sorted index rather than the local shard index, the same input image received the same seed regardless of whether it was processed in the initial run, a resumed run, or a different GPU shard.

\subsection{Brain MRI}

For the Figshare brain MRI test set, target-suppressed images were generated using a conditional DDPM implemented in \texttt{DDPM/sample\_figshare\_dilated\_sharded.py} and executed through \texttt{DDPM/run\_figshare\_test\_dilation9\_3gpu.sh}. The generation set contained 3,064 image--mask pairs. Sampling used the exponential-moving-average weights stored in \texttt{checkpoint/step\_122000.pt}, rather than the raw model parameters. The selected checkpoint corresponded to global step 122,000, and the EMA decay was 0.9999.

The conditional UNet received a three-channel input composed of the masked baseline image, the lesion mask, and a Gaussian-noise channel. It predicted a single-channel diffusion noise term $\epsilon$. The model used 128 base channels, two residual blocks per resolution, channel multipliers of $(1, 1, 2, 2, 4, 4)$, and an attention resolution corresponding to a downsampling factor of 16. Dropout was set to zero, learned variance prediction was disabled, and the model contained approximately 113.67 million parameters.

The diffusion process consisted of 1,000 steps with a linear beta schedule increasing from 0.0001 to 0.02. The model was trained to predict $\epsilon$ using an MSE objective with fixed-small reverse-process variance and no timestep rescaling. Generation used the full 1,000-step DDPM reverse process implemented by \texttt{p\_sample\_loop\_known}; DDIM sampling was not used. Denoised predictions were clipped by setting \texttt{clip\_denoised=True}.

The checkpoint had been trained with a batch size of 8, a learning rate of $1 \times 10^{-4}$, automatic mixed precision, and a random seed of 42. Although the configured training horizon was 2,850,000 steps, the samples used in the present study were generated from the checkpoint saved at step 122,000. During training, the \texttt{synthetic\_healthy} mask mode generated artificial holes by translating, rotating, and flipping tumor-shaped masks and placing them in non-tumor brain regions. The actual tumor was hidden from the conditioning context by setting \texttt{hide\_tumor\_in\_context=True}. The tumor-region loss weight was set to 0, whereas the synthetic-hole and surrounding-context loss weights were set to 1.0 and 0.05, respectively.

At inference time, each grayscale MRI image was first clipped to its 0.1st and 99.9th intensity percentiles and then min--max normalized to the interval $[0,1]$. A region of interest was defined around the bounding box of the actual tumor mask. Its side length was determined as

\[
\max\!\left(
224,\;
h_{\mathrm{bbox}}+32,\;
w_{\mathrm{bbox}}+32
\right),
\]

corresponding to a 16-pixel context margin around each side of the bounding box. The ROI was resized to $224 \times 224$ using bilinear interpolation, and the mask was resized using nearest-neighbor interpolation.

The final application mask was dilated using \texttt{dilation\_size=9}, corresponding to a $9 \times 9$ maximum filter and an expansion of up to four pixels in each direction. This definition differs from the breast-ultrasound setting: \texttt{dilation\_px=9} in the ultrasound pipeline denotes a $19 \times 19$ filter with a nine-pixel radius, whereas \texttt{dilation\_size=9} in the MRI pipeline denotes a $9 \times 9$ filter with a four-pixel radius.

The generated prediction was smoothed using a Pillow Gaussian blur with a radius of 1.075 and resized back to the original ROI size using bilinear interpolation. Two output forms were produced. The raw output replaced the entire ROI with the generated prediction, whereas the blended output replaced pixels only inside the dilated lesion mask using a hard binary boundary. The 3,064 target-suppressed images included in the AUROC and FPR evaluation CSV were taken from the blended outputs.

For an image at sorted dataset index $i$, the random seed was set to $42+i$. The data were partitioned into three shards, with shard $r$ processing indices $r, r+3, r+6, \ldots$. According to the generation logs, shard 0 was processed on an NVIDIA RTX PRO 5000 GPU, whereas shards 1 and 2 were processed on NVIDIA RTX A6000 GPUs.

\subsection{Chest CT}

For the chest CT dataset, $I^{-k}$ was constructed by directly corrupting the target anatomical region with additive Gaussian noise. No generative inpainting model, pixel replacement, or alpha blending was used. The original $512 \times 512$ grayscale JPEG image was first replicated across three channels to obtain an RGB representation. For each pixel $p$ inside the mask of target structure $k$ and each RGB channel $c$, noise was independently sampled as

\[
\epsilon_{p,c} \sim \mathcal{N}\!\left(0,50^2\right),
\]

and the target-suppressed pixel was computed as

\[
I^{-k}_{p,c}
=
\operatorname{clip}\!\left(
I_{p,c}+\epsilon_{p,c},
0,255
\right).
\]

Pixels outside the selected target mask were not explicitly perturbed. The output was converted to 8-bit unsigned integer format after clipping.

The target regions were specified by color-coded masks: blue for the lungs, green for the heart, and red for the bronchi. No mask dilation, erosion, feathering, or boundary smoothing was applied. Consequently, the transformation used an effectively hard mask boundary. A base seed of 20250728 was used. For test image index $i$, the seeds were defined as

\[
\begin{aligned}
s_{i,\mathrm{bronchus}} &= 20250728 + 10i,\\
s_{i,\mathrm{heart}} &= 20250728 + 10i + 1,\\
s_{i,\mathrm{lung}} &= 20250728 + 10i + 2.
\end{aligned}
\]

Thus, independently sampled noise was assigned to every image--structure combination. The resulting images were stored as RGB JPEG files with quality 95 and 4:2:0 chroma subsampling.

We additionally verified the implemented transformation using 100 paired original and target-suppressed images. Within the target masks, the residuals had an empirical mean close to zero, a standard deviation close to 50, and a variance close to 2,500, consistent with the intended additive Gaussian model. The regression slopes between the original and modified pixel intensities ranged from approximately 0.85 to 0.97. The deviation from the ideal additive-model slope of 1 was consistent with clipping at the valid 8-bit intensity limits, particularly in low-intensity regions.

Neighboring residual pixels exhibited weak correlations of approximately 0.05--0.15, supporting the use of spatially independent white noise rather than blurred or spatially correlated noise. Correlations between noise patterns from different images were approximately zero, indicating that noise patterns were not reused. Although noise was independently sampled for each RGB channel before saving, the decoded JPEG files exhibited inter-channel residual correlations of approximately 0.80--0.84 because of chroma subsampling and compression.

Changes outside the target masks were minimal. Using an absolute intensity-difference threshold of 20, the proportions of changed outside-mask pixels were approximately 0.3\%, 0.5\%, and 2.5\% for the lung-, heart-, and bronchus-corrupted images, respectively. These changes were concentrated within approximately four pixels of mask boundaries and were attributable to JPEG block-transform and chroma-subsampling artifacts rather than to intentional noise application.

\subsection{Natural images}

The current documentation identifies the released COD10K and MAS3K datasets but does not specify the procedure used to transform an original image $I$ into its no-object counterpart $I^{-k}$. A reproducible manuscript description therefore cannot yet be written for this domain from the available record alone. The final appendix should report the editing or inpainting method, target-mask preprocessing, spatial resizing, replacement or blending rule, random seed, and output encoding used to generate the COD10K and MAS3K no-object images.

\section{Construction of Positive and Target-Suppressed Captions $T$ and $T^{-k}$}

Let $T$ denote the positive or base caption associated with the original image $I$, and let $T^{-k}$ denote the corresponding caption after suppressing the semantic information associated with target $k$. The caption-generation strategy was selected according to the available labels and the intended evaluation protocol. Most classwise evaluation captions were deterministic class-level templates rather than image-specific free-form descriptions.

\subsection{Breast ultrasonography}

For the UDIAT dataset, class labels were derived from the original Benign and Malignant directories and from the directory containing the generated lesion-suppressed images. These sources were mapped to the labels benign, malignant, and normal, respectively. The evaluation table contained 326 rows: 163 original tumor images and 163 generated target-suppressed images.

Each class was assigned one fixed caption:

\begin{itemize}
  \item benign $\rightarrow$ ``There is a benign tumor.''
  \item malignant $\rightarrow$ ``There is malignant tumor.''
  \item normal $\rightarrow$ ``There is no visible tumor.''
\end{itemize}

These strings were directly defined by the \texttt{CLASSES} constant in \texttt{a1\_inference/build\_udiat\_classwise\_eval\_csv.py}. No language model was queried during construction of this classwise evaluation CSV. The current documentation does not separately specify the caption-construction procedure used for BUS-CoT.

The prompt used to specify the fixed template was:

\begin{quote}
Generate fixed English captions for a three-class breast-ultrasound evaluation. Return exactly three lines only, in this order: benign, malignant, normal. Do not print labels, numbering, quotation marks, or explanations. For benign and malignant, use the literal form ``There is [class] tumor.'', replacing [class] only with benign or malignant. Do not add, remove, or grammar-correct any article. For normal, output ``There is no visible tumor.''
\end{quote}

\subsection{Brain MRI}

For the Figshare brain MRI evaluation set, the tumor type of each of the 3,064 original images was obtained from the \texttt{label\_name} field in \texttt{data/Figshare/test\_figshare3064/figshare\_tumor\_type\_metadata0.json}. The 3,064 DDPM-generated blended images with matching identifiers were assigned the class normal, resulting in a total of 6,128 evaluation rows.

One fixed caption was assigned to each class:

\begin{itemize}
  \item meningioma $\rightarrow$ ``There is a meningioma tumor.''
  \item glioma $\rightarrow$ ``There is a glioma tumor.''
  \item pituitary $\rightarrow$ ``There is a pituitary tumor.''
  \item normal $\rightarrow$ ``There is no visible tumor.''
\end{itemize}

The four strings were defined in the \texttt{CAPTIONS} dictionary of \texttt{a1\_inference/build\_figshare\_classwise\_eval\_csv.py}. Tumor-type metadata were used only to select the corresponding fixed template; no sequence type, tumor location, grade, enhancement pattern, or other image-specific finding was incorporated into the caption.

The prompt used to specify the fixed template was:

\begin{quote}
Generate fixed English captions for a four-class brain-MRI evaluation. Return exactly four lines only, in this order: meningioma, glioma, pituitary, normal. Do not print labels, numbering, quotation marks, or explanations. For each tumor class, use exactly ``There is a [class] tumor.'' and replace [class] only with the supplied class name. For normal, output exactly ``There is no visible tumor.''
\end{quote}

\subsection{Chest CT}

The Chest CT AUROC/FPR evaluation used a deterministic four-class caption set, defined directly in the \texttt{CLASS\_SPECS} constant of \texttt{a1\_inference/build\_chest\_ct\_classwise\_eval\_csv.py}. The evaluation CSV contained 800 rows: 200 images for each of bronchus noise, heart noise, lung noise, and normal. The fixed captions were:

\begin{itemize}
  \item bronchus noise $\rightarrow$ ``Chest CT image showing preserved cardiac contour and normal-appearing bilateral lung parenchyma.''
  \item heart noise $\rightarrow$ ``Chest CT image showing preserved bilateral lung aeration and patent central airways.''
  \item lung noise $\rightarrow$ ``Chest CT image showing preserved cardiac contour and patent central tracheobronchial structures.''
  \item normal $\rightarrow$ ``Chest CT image showing preserved cardiac contour, preserved bilateral lung aeration, and patent central airways.''
\end{itemize}

The normal caption serves as the all-structure baseline $T$. For each regional-noise class, the target-suppressed caption $T^{-k}$ describes only the two structures left uncorrupted by Gaussian noise and omits the corrupted target structure. The same deterministic captioning rule was used when constructing the Chest CT DPO training data; no separate prompt or caption-generation procedure was introduced for the AUROC/FPR evaluation. No per-image caption generation, random sampling, or additional language-model inference was used in this classwise evaluation.

\subsection{Natural images}

For COD10K, the original class labels were obtained from the 69 camouflage subclasses provided in the test metadata. A further no camouflaged animal class was assigned to the generated no-object images, resulting in 70 evaluation classes. For MAS3K, the 33 original object classes were supplemented with the same no-object class, resulting in 34 evaluation classes.

For an original image containing an annotated camouflaged object, the caption was deterministically constructed as ``There is an [class name] present in the scene.'', where [class name] was replaced by the corresponding dataset class. The article \emph{an} was intentionally retained for every class and was not grammatically adjusted according to the initial sound of the class name. For a generated no-object image, the fixed target-suppressed caption was ``There is no camouflaged animal present in the scene.''

No habitat, color, size, background description, or image-specific attribute was appended. The same deterministic template rule was used for the COD10K and MAS3K DPO training-data construction and was reused for classwise AUROC/FPR evaluation; no separate prompt or caption-generation procedure was introduced for either dataset.

The COD10K evaluation file contained 2,026 original images and 2,026 no-object images, yielding 4,052 rows across 70 classes. The corresponding file was \texttt{data/COD10K/test/cod10k\_original2026\_no\_object2026\_subclass70\_scene\_caption\_pairs\_paths\_fixed.csv}.

The MAS3K evaluation file contained 582 original images and 582 no-object images, yielding 1,164 rows across 34 classes. The corresponding file was \texttt{data/MAS3K/mas3k\_original582\_no\_object582\_classwise\_scene\_caption\_pairs\_paths\_fixed.csv}.
